% sup3_02_S13: Supplement part 3, the properties of the cube split (S12, printed) used below, and step S13: the constant field B_0.

\subsection{The cube split: the three properties used}\label{sup3:ssec:S12}

\cite[p.~433]{Ba3}: ``We take a cube $\square_1$ of size $r(L^k\eps)$ containing the above cube in the center, and next we take a cube $\square$ of size $3r(L^k\eps)$ and with $\square_1$ in the center. We assume that $\square_1$, $\square$ are sums of big blocks of the unit lattice. Now we replace each scalar field propagator $G_k(\Omega,\tilde B)$ [or $G_k(\Omega_2,\tilde B)$] in the expression by $G_k(\square,\tilde B)$, more exactly we represent $G_k(\Omega,\tilde B)=G_k(\square,\tilde B)+\delta G_k(\Omega,\square,\tilde B)$, and we treat $\delta G_k(\Omega,\square,\tilde B)$ as an external field. Thus again we drop the lines corresponding to these operators and decompose the obtained set of graphs into renormalized classes.'' This step is printed. The main paper performs it with the constant side $r^\flat$ in place of $r(t)$ (the containment $\square\subset\Omega$ fails at $L=5$ for the printed side) and proves the bound on $\delta G_k(\Omega,\square,\tilde B)$ there. The steps S13--S15 use only the following three properties, which hold for $r_\square=r^\flat$ under the main paper's cube step and for $r_\square=r(t)$ wherever the printed step applies.
\begin{enumerate}[label=\textup{(P-\roman*)},nosep,leftmargin=*]
\item $\square_1\subset\square$ are unions of large blocks, obtained by rounding outward the cubes of sides $r_\square$ and $3r_\square$ concentric with the class-component cube $Q_0$ of \cref{sup3:cor:S114}(c). So $\square_1$ contains the cube of side $r_\square$ concentric with $Q_0$, $\square$ contains the cube of side $3r_\square$ concentric with $Q_0$, and $\operatorname{side}(\square)\le3r_\square+2M$.
\item $\square\subset\Omega$ (resp.\ $\square\subset\Omega_2\subset\Omega$), so $\tilde B$ is defined, and satisfies the hypotheses of \cite[p.~412]{Ba3}, on every bond of $\square$; and $\square$ does not wrap around the torus.
\item After the split every scalar line is $G_k(\square,\tilde B)$; $\delta G_k(\Omega,\square,\tilde B)$, $\delta G_k(\Omega,\Omega_2,\tilde B)$ and the kernels \cite[(1.16)]{Ba3} are external fields.
\end{enumerate}

\subsection{S13: the constant field $B_0$}\label{sup3:ssec:S13}

\cite[p.~433]{Ba3}: ``The field $\tilde B$ is regular in the sense that $\norm{\partial^\eta_\mu\tilde B}_{0,\alpha}\le O(p(L^k\eps))$. Hence there exists a constant field $B_0$ such that $|\tilde B-B_0|\le O(r(L^k\eps)p(L^k\eps))\le O(p(L^k\eps)^2)$ on the cube $\square$.'' The consumed form: a constant vector field $B_0$ with $|\tilde B-B_0|\le Cr_\square p(t)$ on the bonds of $\square$, such that $\tilde B'\defeq\tilde B-B_0$ has the derivative and H\"older bounds of $\tilde B$ and is an external vector field of norm \cite[(1.32)]{Ba3} $O((r_\square+1)p(t))$ on $\square$: $O(p(t))$ at the constant side $r^\flat$.

\begin{lemma}[the constant field]\label{sup3:lem:S13}
Assume (P-i), (P-ii), and let $c$ be a site of $\square$ with $c+\eta e_\mu\in\square$ for every $\mu$ and $|x-c|_\infty\le\frac12\operatorname{side}(\square)+\eta$ for all $x\in\square$ (a site next to the centre). Put $B_{0,\mu}\defeq\tilde B_\mu(c)$ for every $\mu$ (a constant field) and $\tilde B'\defeq\tilde B-B_0$ on the bonds of $\square$. Then for every bond $(x,x+\eta e_\mu)\subset\square$ and every contour $\Gamma\subset\square$:
\begin{enumerate}[label=\textup{(\alph*)},nosep,leftmargin=*]
\item $|\tilde B'_\mu(x)|\le c_1p(t)\,|x-c|_1\le C_{13}\,p(t)$, with $C_{13}\defeq c_1d\big(\tfrac32r_\square+M+1\big)$;
\item $\partial^\eta\tilde B'=\partial^\eta\tilde B$ on $\square$, so $|\partial^\eta_\nu\tilde B'_\mu|\le c_1p(t)$, and the H\"older quotients of $\partial^\eta\tilde B'$ are those of $\partial^\eta\tilde B$;
\item $|\tilde B'(\Gamma)|\le|\Gamma|_\eta\,C_{13}\,p(t)$, with $|\Gamma|_\eta$ the length of $\Gamma$ ($\le d$ for the contours $\Gamma^{(k)}_{y,x}$);
\item the norm \cite[(1.32)]{Ba3} of $\tilde B'$ as an external vector field (``with $\tilde B=0$'', \cite[p.~420]{Ba3}) over $\square$ satisfies $\norm{\tilde B'}_{1,\alpha_0;\square}\le(C_{13}+3c_1)p(t)$ wherever the H\"older clause on $\tilde B$ holds, i.e.\ on pairs of points at depth $\ge R_{\mathrm{Ba}}+2$ in $\Omega$;
\item at $r_\square=r^\flat$, $C_{13}$ is a constant and $\tilde B'=O(p(t))$, the size class of $\tilde A$; at the printed $r_\square=r(t)$, $|\tilde B'|\le c_1d(\tfrac32R(1+\ln t^{-1})^{\varrho}+M+1)p(t)$, which is $\le C(R/b_0+1)p(t)^2$ uniformly in $t\le1$ if and only if $\varrho\le\bar p$.
\end{enumerate}
\end{lemma}

\begin{proof}
\emph{A path inside $\square$.} Fix a bond $(x,x+\eta e_\mu)\subset\square$. Go from $c$ to $x$ by changing one coordinate at a time, $\nu=1,\dots,d$, in steps of $\pm\eta$; every intermediate site $z$ has $z_\nu$ between $c_\nu$ and $x_\nu$ for the coordinates already moved and equal to $c_\nu$ or $x_\nu$ otherwise. $\square$ is a product of intervals of sites (a cube that does not wrap, (P-ii)); both $c+\eta e_\mu$ and $x+\eta e_\mu$ lie in $\square$, so $\max(c_\mu,x_\mu)+\eta$ is at most the last site coordinate of $\square$ in direction $\mu$, and $z+\eta e_\mu\in\square$ for every intermediate $z$. Hence every pair of bonds $(z,z+\eta e_\mu)$, $(z+\eta e_\nu,z+\eta e_\nu+\eta e_\mu)$ used below lies in $\square\subset\Omega$, where the hypothesis on $\partial^\eta\tilde B$ holds.

\emph{(a) Telescoping.} $\tilde B_\mu(x)-\tilde B_\mu(c)=\sum_{\text{steps}}\pm\eta\,\partial^\eta_\nu\tilde B_\mu(z)$ over $|x-c|_1/\eta$ steps, each bounded by $\eta c_1p(t)$. So $|\tilde B'_\mu(x)|\le c_1p(t)|x-c|_1\le c_1p(t)\,d(\tfrac12\operatorname{side}(\square)+\eta)\le C_{13}p(t)$, using $\tfrac12\operatorname{side}(\square)\le\tfrac32r_\square+M$ (P-i) and $\eta\le1$.

\emph{(b)} $B_0$ is constant, so $\partial^\eta B_0=0$.

\emph{(c)} $\tilde B'(\Gamma)=\sum_{b\subset\Gamma}\eta\tilde B'_b$ has $|\Gamma|_\eta/\eta$ terms, each bounded by $\eta C_{13}p(t)$ by (a).

\emph{(d)} The norm is the sum of the sup (a), the derivative bound (b) and the H\"older bound: for $|x-x'|\le1$ the quotient is $\le c_1p(t)$ by (b) and the hypothesis; for $|x-x'|>1$ it is $\le2\sup|\partial^\eta\tilde B'|\le2c_1p(t)$.

\emph{(e)} Insert $r_\square$. At the printed cube, $r(t)p(t)=(R/b_0)(1+\ln t^{-1})^{\varrho-\bar p}p(t)^2$, which is $O(p(t)^2)$ uniformly in $t\le1$ exactly when $\varrho\le\bar p$.
\end{proof}

\begin{remark}\label{sup3:rem:S13}
(a) \emph{The inequality ``$O(rp)\le O(p^2)$''.} \cite[(2.3), (2.7)]{Ba2} fix only $\bar p>2$ and $\varrho>1$; the printed inequality needs $\varrho\le\bar p$, which holds for the values used in the main paper ($\varrho=2$, $\bar p=6$) and is not needed at the constant side, where $\tilde B'=O(p)$ by (e).

(b) \emph{Where the H\"older clause is used.} The norm (d) is needed at the vertices of the $\tilde B'$-expansion (\cref{sup3:ssec:S14b}). At the vertices of the components containing an old vertex the depth is $\ge r_\square$ (\cref{sup3:lem:S15}), and the clause holds there when $r^\flat\ge R_{\mathrm{Ba}}+2$. At the vertices of the chains without an old vertex (S16, the fourth part), which may lie anywhere in $\square$, a point $z\in\square$ has depth $\ge\rho-D_\square$, with $\rho$ the depth of the old vertices and $D_\square\defeq\tfrac32r_\square+2M+l_N$; the main paper's cube step places $Q_0$ at depth $>D^\flat=\tfrac32r^\flat+2M+l^\sharp$, so $\operatorname{depth}(z)>l^\sharp-l_N=4V^\sharp_{\mathrm{lab}}$ (\cref{sup3:ssec:S15}). The parameter row $R\ge R_{\mathrm{flat}}+R_{\mathrm{Ba}}+2$ of the main paper ($R_{\mathrm{flat}}$ the depth parameter of its cube step) makes every point of $\square$ lie at depth $\ge R_{\mathrm{Ba}}+2$ in $\Omega$, so (d) holds on all of $\square$; it is a lower bound on $R$ of Ba{\l}aban's kind, fixed after $R_{\mathrm{Ba}}$ and $l^\sharp$. The chain step of the fourth part in fact uses only $\sup|\tilde B'|$ and $\sup|\partial^\eta\tilde B'|$ at the chain vertices, so the row is a convenience there.

(c) \emph{On a torus} the same proof works for a cube $\square$ that does not wrap (P-ii), so that the coordinate path is defined.

(d) Nothing in the lemma depends on $k$ except through $c_1$, whose $k$-uniformity is the input named in the conventions.
\end{remark}
